\subsection*{首次提交与查询}

现在把前面的选择收成一组可检查的接口规则。购买动作采用 \code{purchase.create.v1}。服务从认证上下文获得账户 1，客户端在首次发送前保存 K1、Q7 和输入。K1 是示意值，实际键应满足公布的字符和长度要求。每次 HTTP 尝试使用新的请求 ID。

\begin{BookCode}
POST /v1/purchases HTTP/1.1
Authorization: Bearer <access_token>
Idempotency-Key: K1
Request-Id: R1
Operation-Permit: <operation_permit>
Content-Type: application/json

{"product_id":"paper-001","quantity":1,
 "currency":"CNY","quote_id":"Q7"}
\end{BookCode}

\code{Authorization} 携带访问凭据，服务校验后得到调用者及其可操作账户。\code{Request-Id} 用于本次通信诊断，\code{Content-Type} 声明正文按 JSON 解析。

\code{Operation-Permit} 是服务签发、可核验真伪的操作许可，把账户、动作、登记会话 E1、K1 和固定截止时间联系起来。登记会话是一组共同受固定期限约束的操作身份，作用与登录会话不同。操作许可提供引用身份的资格，购买访问权仍由账户认证与授权检查决定；它也没有提前指定哪份商品输入一定成交。

服务核验许可，以账户、动作、E1 和 K1 查找原记录。存在记录时核对规范化输入，已完成则返回保存的结果；新身份进入购买处理。以下是一项成功购买的业务结果示例，\code{retry\_until} 表示这项身份允许恢复的截止时刻，末尾 Z 表示采用 UTC，即协调世界时。

\begin{BookCode}
{
  "idempotency_key": "K1",
  "purchase_id": "P1",
  "status": "succeeded",
  "order_id": "O1",
  "debit_id": "D1",
  "product_id": "paper-001",
  "currency": "CNY",
  "amount_minor": 3000,
  "balance_after_minor": 7000,
  "quote_id": "Q7",
  "retry_until": "2026-10-18T10:00:00Z"
}
\end{BookCode}

同键同参的成功重试返回相同购买、订单、扣款和成交快照。\code{balance\_after\_minor} 始终表示 P1 扣款后的余额。客户端要展示当前可用余额时，再访问账户查询入口。本次请求 ID 或重放标志可以随响应变化，它们与保存的购买结果分开。

查询入口为 \code{GET /v1/purchase-operations/K1}，请求继续携带能确定 E1 的有效许可，并在已认证账户与 \code{purchase.create.v1} 空间内解释 K1。有效期内，200 返回已保存终态；404 表示当前读取路径尚未找到该身份的记录。404 仍允许原 POST 在途中，客户端保留 K1，按同一输入恢复。查询自身没有创建订单或撤销原请求的作用。

\begin{table}[htbp]
\centering\small
\begin{tabularx}{\linewidth}{@{}p{4.1cm}XX@{}}
\toprule
HTTP 状态与业务码 & 接口含义 & 客户端继续方式\\
\midrule
201 \code{succeeded} & 首次提交形成成功终态 & 保存 P1 和原结果\\
409 \code{rejected} & 首次提交形成余额不足终态 & 保存原拒绝\\
200 \code{succeeded/rejected} & 找到原身份的已保存终态 & 按原结果结束本地操作\\
409 \code{conflict} & 同身份绑定了不同输入 & 核对旧记录，提示冲突\\
400 \code{invalid\_input} & 语法解析失败，本次未受理 & 修正语法后按规则提交\\
422 \code{invalid\_input} & 字段校验失败，本次未受理 & 修正字段后按规则提交\\
503 \code{unavailable} & 暂时无法完成处理或确认 & 保留 K1，恢复原操作\\
410 \code{operation\_expired} & 身份的恢复资格已结束 & 查询业务事实，另行确认\\
\bottomrule
\end{tabularx}
\caption{状态码与业务码共同说明本接口的处理结果。}
\end{table}

现有四表事务中 \code{pending} 只存在于事务内部，外部普通查询看不到未提交行；因此查询只提供已保存终态或无记录这两类证据。若要对外提供处理中状态，必须另行设计可提交的受理记录及其执行、接管和恢复规则。

\subsection*{把有效范围写进身份规则}

逐笔结果若最终需要删除，服务就要拒绝已经失效的旧身份。否则，曾经执行后被清理的 K1 与从未出现的新键都呈现为无记录，迟到的旧请求可能再次扣款。RIFL 的租约与 Ongaro 的注册会话都为旧调用规定了可检查的资格边界；其中找不到会话时自动注册新会话，会使旧命令重新取得执行机会\cite{rifl,ongaro}。

这里采用固定期限的登记会话。服务先持久登记 E1 所属账户、动作、截止时间和 \code{active} 状态，入口只接受已登记、仍有效的会话。客户端选定并持久保存 E1 后，与 K1 一起使用。完整身份由此为
\[
 S=(\mathrm{account\_id},\,\mathrm{action},\,\mathrm{session\_id},\,
       \mathrm{idempotency\_key}).
\]
\code{session\_id} 在例子中取 E1。同一会话中的新购买仍生成新 K；同一购买的所有尝试则始终保留原 E1、K1 与确认输入。已经发出购买以后改带 E2，会形成另一个完整身份，原购买仍待查明。

本书采用以下关闭边界：E1 的截止时间固定，关闭决定一旦持久保存就不可撤回；旧 E1 不续期、不重新注册，也不分配给新的有效会话。入口对关闭或缺失的会话都拒绝旧请求，重启或切换后继续保留这项判断依据。清理还要保护关闭以前已经受理的在途执行，使其身份与结果证据一直保留到安全终结。关闭须先阻止新的受理，已有执行提交或回滚后，才可清理相关逐笔记录。

\begin{figure}[htbp]
\centering
\begin{tikzpicture}[x=1cm,y=1cm]
\node[booknode,text width=3.6cm] (c) at (1.8,0) {可核验会话 E1\\固定截止 10 月 18 日 10:00\\操作使用 (E1, K1)};
\node[booknode,text width=3.6cm] (r) at (6.5,0) {截止前的同身份恢复\\找到 P1，核对输入\\返回原结果};
\node[bookresult,text width=3.6cm] (e) at (11.2,0) {持久关闭后的旧 E1\\关闭或缺失均拒绝\\返回 410};
\draw[bookarrow] (c)--(r);
\draw[bookalert] (r)--(e);
\node[booknote,text width=12cm] at (6.5,-1.6)
 {先持久关闭会话，再等待在途执行安全终结，最后清理逐笔结果\\关闭不可撤回；缺失拒绝；旧 E1 不重新注册};
\node[booknode,text width=10.5cm] at (6.5,-2.9)
 {完整身份：账户、动作、登记会话 E1、幂等键 K1\\新 E2 属于另一身份，切换会话不能查清旧 E1 的购买};
\end{tikzpicture}
\caption{完整身份在有效会话内使用，关闭后旧身份不再取得受理资格。}
\label{fig:ch03-session-expiry}
\end{figure}


例如 E1 在 10 月 18 日 10:00 截止。此前完成的 P1 可以在有效期内按原身份恢复；关闭以后，旧提交入口返回 \code{410 operation\_expired}。410 只说明恢复资格结束，P1 仍可能已经成功，客户端应查询账户内的购买与订单事实，必要时转入受控核查。用户明确确认再买一份，才为另一项购买取得新身份。

第一部分在会话有效、逐笔记录保留的范围内实现购买与恢复。许可的防篡改、判期时间、时钟回拨以及关闭与在途执行如何协调，需要完整的回收协议，第十二章将继续推导这些条件。

\subsection*{服务要保存什么证据}

服务重新启动后，内存中曾经见过 K1 的事实已经消失。请求再次到达时，应从持久记录中得到 K1 与 P1 的关联。副本切换后若读到尚未包含 P1 的旧副本，也可能误判为新操作，查找身份的读取路径必须与后续受理机制一起设计。

一条已完成身份记录至少包含经过验证的账户与动作、E1 与 K1、规范化版本及输入、首次受理时间、状态、P1、原结果与恢复截止时间。成功结果关联 O1、D1 及成交快照；拒绝结果保存明确业务原因。Q7 的内容或可长期解释它的快照也应保留。只保存一个处理过标记，会失去原购买编号和原返回值，客户端仍然无法恢复完整结果。

四张业务表的关系保持不变。\code{accounts} 保存余额，\code{purchases} 表示业务处理，\code{orders} 保存成交，\code{debits} 保存资金变动。身份映射可以作为 \code{purchases} 的受理字段，也可以单独建立操作记录表。单独建表时，需要把成功映射、终态和订单扣款纳入同一次原子提交；否则可能先扣款而没有映射，或映射说成功而订单尚未提交。

\begin{figure}[htbp]
\centering
\begin{tikzpicture}[x=1cm,y=1cm]
\node[booknode,text width=5.3cm] (s) at (2.8,0) {身份与输入证据\\账户、动作、E1、K1\\规范化版本、输入、恢复期限};
\node[booknode,text width=5.3cm] (p) at (10,0) {购买记录 P1\\状态、首次受理时间\\稳定原结果与成交快照};
\draw[bookarrow] (s)--node[above,font=\footnotesize]{定位}(p);
\node[booknode,text width=3.6cm] (o) at (7.4,-1.8) {成功订单 O1\\引用 P1};
\node[booknode,text width=3.6cm] (d) at (12.1,-1.8) {扣款流水 D1\\金额、当次扣后余额};
\draw[bookarrow] (p)--(o);
\draw[bookarrow] (p)--(d);
\begin{scope}[on background layer]
\node[draw=Accent!65,dashed,rounded corners=4pt,fit=(s)(p)(o)(d),inner sep=10pt] {};
\end{scope}
\node[booknote,text width=12cm] at (7,-3.25)
 {共同提交范围：身份绑定、业务变化、终态与原结果\\拒绝终态也保留身份、输入及拒绝原因};
\end{tikzpicture}
\caption{身份事实需要与原购买的记录关系一起保留，恢复才能返回正确的业务结果。}
\label{fig:ch03-evidence}
\end{figure}


\FloatBarrier
为了观察身份规则的作用，可以写出一个最小串行流程：
\begin{BookCode}
验证账户、动作、会话有效性与请求格式
按适用的规则规范化输入，得到版本和指纹
查找完整操作身份 S
若已有终态：核对输入并返回原结果
若是新操作：验证报价，执行购买事务
在业务提交中保存身份、终态和原结果
提交确认后发送响应
\end{BookCode}

适用的规范化规则包括原记录保存的版本。实际实现若需先读取版本，再解释本次输入，应按这个依赖安排顺序。这个流程先按同身份请求依次进入推演。身份、原结果和业务变化若纳入同一事务，崩溃后只有共同提交或共同回滚两种恢复结果。

并发服务器上，R1 与 R2 却可能同时查到无记录，各自提交购买。新建的 P、O、D 编号不同，各自的唯一约束仍允许它们通过。服务也可能发送提交后失去数据库连接。下一步需要在共享数据库中保护完整操作身份，让重试恢复与业务效果共同提交的原结果。
\FloatBarrier
